\documentclass[11pt,a4paper]{article}
\usepackage[T1]{fontenc}
\usepackage[utf8]{inputenc}
\usepackage{lmodern}
\usepackage{amsmath,amssymb,amsthm,mathtools}
\usepackage[a4paper,margin=26mm,headheight=15pt]{geometry}
\usepackage{microtype,booktabs,tabularx,array}
\usepackage{enumitem}
\usepackage{xcolor,fancyhdr}
\usepackage{hyperref}
\usepackage[nameinlink,noabbrev]{cleveref}
\definecolor{inkblue}{RGB}{28,55,81}
\hypersetup{colorlinks=true,linkcolor=inkblue,citecolor=inkblue,urlcolor=inkblue,
 pdftitle={From Analytic Spectral Collapse to Sobolev Spectral Disks},
 pdfauthor={OpenAI ChatGPT; research manuscript prepared for Vladimir Reshetnikov},
 pdfsubject={Exact spectra, Thue--Morse and Stern dual distributions, and sharp logarithmic endpoint asymptotics}}
\pagestyle{fancy}\fancyhf{}
\fancyhead[L]{\small Sobolev spectral disks and endpoint laws}
\fancyhead[R]{\small Research manuscript}
\fancyfoot[C]{\thepage}
\setlength{\parskip}{0.3em}
\setlength{\parindent}{1.15em}
\setlength{\emergencystretch}{2em}
\setlist[enumerate]{itemsep=0.3em,topsep=0.3em}
\newtheorem{theorem}{Theorem}[section]
\newtheorem{lemma}[theorem]{Lemma}
\newtheorem{proposition}[theorem]{Proposition}
\newtheorem{corollary}[theorem]{Corollary}
\theoremstyle{definition}
\newtheorem{definition}[theorem]{Definition}
\newtheorem{question}[theorem]{Question}
\newtheorem{example}[theorem]{Example}
\theoremstyle{remark}
\newtheorem{remark}[theorem]{Remark}
% ed. (2026-09-29): unnumbered environment for editorial notes.
\newtheorem*{ednote}{Editorial note (ProveIt, 2026-09-29)}
\newcommand{\C}{\mathbb C}
\newcommand{\R}{\mathbb R}
\newcommand{\Q}{\mathbb Q}
\newcommand{\Z}{\mathbb Z}
\newcommand{\N}{\mathbb N}
\newcommand{\T}{\mathbb T}
\newcommand{\dd}{\,\mathrm d}
\newcommand{\ii}{\mathrm i}
\newcommand{\Lop}{\mathcal L}
\newcommand{\ess}{\mathrm{ess}}
\newcommand{\Span}{\operatorname{span}}
\newcommand{\ran}{\operatorname{ran}}
\newcommand{\tr}{\operatorname{tr}}
\newcommand{\sinc}{\operatorname{sinc}}
\newcommand{\norm}[1]{\left\lVert#1\right\rVert}
\newcommand{\abs}[1]{\left|#1\right|}
\newcommand{\br}[1]{\langle#1\rangle}
\newcommand{\Disk}[1]{\overline{D}(0,#1)}
\newcommand{\Hs}{H^{s,\beta}(\T)}
\numberwithin{equation}{section}

\title{\textbf{From Analytic Spectral Collapse\\to Sobolev Spectral Disks}\\[0.6em]
\large Exact regularity thresholds and sharp endpoint laws\\
for Rvachev--Thue--Morse transfer operators}
\author{Research manuscript prepared for Vladimir Reshetnikov\\
\small OpenAI ChatGPT}
\date{29 September 2026 (Pacific time)}

\begin{document}
\maketitle
\begin{abstract}
The ProveIt manuscript on spectral collapse and alternating Rvachev RMS
asymptotics asks how its analytic spectrum changes at finite smoothness.
We answer the periodic Sobolev part of that question, including logarithmic
refinements, and the nonperiodic interval problem at every nonnegative integer
Sobolev order. For a trigonometric polynomial mask and an integer dilation,
a positive eigenfunction of the squared-mask transfer operator gives the
exact Fredholm essential spectrum: a closed disk. The complete spectrum is
that disk together with the spectrum of an explicit finite Fourier matrix.
We prove the required positivity hypothesis for every even sine power.

For the dyadic quadratic sine mask the disk radius is
$2^{-s}\sqrt{1+\sqrt{17}}/4$. The eigenvalues $1/2$ and $-1/4$ become isolated
exactly above two regularity thresholds separated by one derivative.
An exact Fourier-cutoff identity links the leading Thue--Morse correlation
functional to the subleading Stern-polynomial functional $B_k(-2)$.
Closed quadratic-energy formulas determine their precise Sobolev and
logarithmic continuity thresholds. They also give two-sided, sharp operator
norms for the scalar moment remainder. At the subleading critical index,
a logarithmic weight of exponent $\beta>1/2$ recovers a uniform two-term
scalar asymptotic even though the subleading eigenvalue remains on the
essential spectral boundary. No corresponding operator-norm little-oh
remainder is possible.

All infinite-dimensional conclusions are proved analytically. An accompanying
standard-library Python program checks finite identities with exact rational
and quadratic-field arithmetic. This is an unrefereed research extension,
not a new Lean formalization or a claim of established global priority.
\end{abstract}
\noindent\textbf{Keywords.} Transfer operator; essential spectrum; Sobolev regularity;
Thue--Morse correlations; Stern polynomials; Riesz product; critical asymptotics.

\clearpage
{\small\setlength{\parskip}{0pt}\tableofcontents}
\clearpage

\section{The repository question and the contribution}
\label{sec:intro}

\subsection{A specific, delimited research target}
The Fabius-function branch of ProveIt contains a Fourier-decay synthesis and
a later manuscript entitled \emph{Spectral Collapse and Alternating RMS
Asymptotics for the Rvachev Up-Function} \cite{repo-spectral,repo-fourier}.
The later manuscript proves, on analytic Hardy disk spaces, that the quadratic
dyadic transfer operator has spectrum $\{0,1/2,-1/4\}$. Its research question
\emph{Spectra at finite smoothness} asks for the spectrum or essential spectral
radius on finite-smoothness spaces, the isolation of the analytic eigenvalues,
and the effect of unmatched endpoint jets. It explicitly warns against
transferring an analytic determinant to such spaces.

This paper answers that question on $H^{s,\beta}(\T)$ for all real $s,\beta$,
and on $H^r(0,1)$ for every integer $r\ge0$. The answer differs fundamentally
from analytic spectral collapse: there is an entire essential spectral disk,
not merely a finite set. The remaining H\"older, $C^r$, and fractional
nonperiodic interval questions are not claimed to be resolved.

The comparison is fixed by the repository commit
\begin{center}\small\texttt{b2aed2d23125006c48d63fc97aef0d64b67becb1}.\end{center}
The exact source path, inspected material, and source identifiers are recorded
in \cref{sec:provenance}. No repository file was modified in preparing this
article. In particular, a result described here as proved is a claim about
the mathematical argument in this article, not a claim that the repository's
Lean development already contains it.

\subsection{What is already known, and what is supplied here}
Transfer operators with sine and cosine weights have an established
literature. Chen and Volkmer study precisely this family on the circle,
including explicit dyadic computations \cite{chen-volkmer}. Disk spectra
for weighted isometries and their connection with transfer operators are
part of a broader theory studied by Bardadyn and Kwa\'sniewski
\cite{bardadyn-kwasniewski}. Thue--Morse correlations, their Riesz-product
origin, and questions about their means are also established objects
\cite{baake-coons}. The recurrence defining Stern polynomials is classical;
we use the normalization of Ulas and Ulas \cite{ulas-ulas}.

The present work does not claim to discover those operators, sequences, or
the general possibility of a disk spectrum. Its contribution is the
self-contained combination of an exact Sobolev essential disk with a
complete finite-matrix classification of the outliers, followed by a sharp
critical-regularity analysis of the two distinguished dyadic coefficient
functionals. The exact Fourier-tail formula yields a scalar asymptotic at an
embedded spectral boundary, together with an optimal logarithmic threshold.
These statements provide a concrete answer to the inspected repository's
finite-smoothness question.

The literature checks establish relevant antecedents but are not an exhaustive
priority investigation. Some retrieved sources were available only through
primary-source abstracts and metadata. Accordingly, the results are presented
as a proved research extension relative to the specified repository target,
not as a certified first occurrence of every specialization in the literature.
% ed. (2026-09-29): uncited overlap with Part I of the Thue--Morse frontier
% deductions and with the predecessor; uncited Lean theorems.
\begin{ednote}
Several ingredients of \cref{sec:cutoff,sec:energies} are already in the
repository, which this article does not cite. Part~I of the volume
\path{thue-morse/Thue_Morse_Frontier_Deductions/} of the Fabius
research-frontier drafts tree (filed 2026-09-05) proves, for the same Riesz product and the same
sequences, the cutoff identity \eqref{eq:cutoff} together with
$a_k=B_k(-2)$, the tail formula of \cref{cor:tail-formula}, the energy
recurrences and closed forms of
\cref{lem:energy-recurrences,prop:closed-energies} (with a trapezoidal
endpoint weight), the exponents $s_0,s_1$, the unweighted ($\beta=0$) case
of \cref{thm:dual-thresholds}, and remainder rates in $H^{-s}$; it credits
the growth rate $1+\sqrt{17}$ to Zaks, Pikovsky and Kurths. The operator
form of \eqref{eq:cutoff} is the finite-mode reduction of the predecessor
\cite{repo-spectral}, whose sequences $\eta,\alpha$ are the present
$\eta,a$ and whose quartic spectrum is \eqref{eq:A22-char}. New here are
the logarithmic refinement, the two-sided remainder laws, the critical
scalar/operator separation and the spectral classification of
\cref{sec:framework,sec:L2,sec:sobolev-proof,sec:sine}. The
eigen-identities \eqref{eq:primal-eigenfunctions} are machine-checked in
\path{Analysis/FabiusFunction/Lean/FabiusFunction/RMSTransferEigenfunctions.lean}
and the recursion for $\eta$ in \path{ThueMorseAutocorrelationLimit.lean}
of the same directory.
\end{ednote}

\subsection{The principal dyadic conclusion}
Let
\begin{equation}
 (\Lop f)(x)=\frac12\left[
 \sin^2\!\left(\frac{\pi x}{2}\right)f(x/2)
 +\cos^2\!\left(\frac{\pi x}{2}\right)f((x+1)/2)\right].
 \label{eq:intro-L}
\end{equation}
Put
\begin{equation}
 \gamma=\frac{\sqrt{1+\sqrt{17}}}{4},\qquad
 s_0=\frac12\log_2\!\left(\frac{1+\sqrt{17}}4\right),\qquad
 s_1=s_0+1.
 \label{eq:thresholds-intro}
\end{equation}
The exact classification proved below is
\begin{equation}
 \sigma_{\ess}(\Lop;H^{s,\beta})=\Disk{\gamma2^{-s}},\qquad
 \sigma(\Lop;H^{s,\beta})=
 \Disk{\gamma2^{-s}}\cup\{1/2,-1/4\}.
 \label{eq:intro-spectrum}
\end{equation}
Here essential spectrum means \emph{failure of Fredholmness}. The real
logarithmic exponent $\beta$ does not change this spectral set. It does,
however, decide whether the leading and subleading coefficient functionals
are continuous at $s=s_0$ and $s=s_1$.

For the sine-product moments
\begin{equation}
 W_n(x)=\prod_{j=0}^{n-1}\sin^2(\pi2^j x),\qquad
 I_n(f)=\int_0^1 f(x)W_n(x)\dd x,
 \label{eq:intro-moments}
\end{equation}
we construct functionals $\mu,D_+,D_-$ such that
\begin{equation}
 I_n(f)=2^{-n}\mu(f)
  +\frac{(-1/4)^n}{3}(D_++D_-)(f)+R_n(f).
 \label{eq:intro-two-term}
\end{equation}
At the critical subleading regularity,
\begin{equation}
 \norm{R_n}_{(H^{s_1,\beta})^*}
 \asymp_\beta 4^{-n}(1+n)^{1/2-\beta}
 \quad(\beta>1/2).
 \label{eq:intro-endpoint}
\end{equation}
Thus a uniform scalar remainder $o(4^{-n})$ survives at the essential
spectral boundary, while an operator-norm remainder $o(4^{-n})$ cannot.
This distinction is one of the main reasons to study both the spectrum and
the actual moment functional rather than substitute one question for the other.

\section{Spaces, masks, and the classification theorem}
\label{sec:framework}

\subsection{Fourier conventions and logarithmic Sobolev spaces}
Write $\T=\R/\Z$, $e_k(x)=\exp(2\pi\ii kx)$, and
$\widehat f(k)=\int_0^1f(x)e^{-2\pi\ii kx}\dd x$ for integrable $f$.
For real $s,\beta$, define $H^{s,\beta}(\T)$ as the Hilbert completion of
trigonometric polynomials in the norm
\begin{equation}
 \norm f_{s,\beta}^2=
 \sum_{k\in\Z}\br{k}^{2s}\bigl[\log(e+|k|)\bigr]^{2\beta}
       |\widehat f(k)|^2,\qquad \br{k}=(1+k^2)^{1/2}.
 \label{eq:Hsb}
\end{equation}
The notation $H^s$ means $H^{s,0}$. For negative $s$, the completion is
naturally a space of periodic distributions. Every formula initially stated
on trigonometric polynomials extends by boundedness when the relevant norm
estimate is established. Pairings with distributions are complex-linear in
$f$; they need not be Hilbert inner products.

Let $b\ge2$ be an integer and let
\begin{equation}
 w(x)=\sum_{\ell=-D}^{D}w_\ell e_\ell(x)\ne0,
 \qquad
 (\Lop_{b,w}f)(x)=\frac1b\sum_{a=0}^{b-1}
       (wf)\!\left(\frac{x+a}{b}\right).
 \label{eq:general-L}
\end{equation}
The Fourier rule is
\begin{equation}
 \widehat{\Lop_{b,w}f}(k)
   =\sum_{\ell=-D}^{D}w_\ell\widehat f(bk-\ell),
 \qquad
 \Lop_{b,w}e_j=\sum_k w_{bk-j}e_k.
 \label{eq:Fourier-rule}
\end{equation}
In particular, these operators are bounded on every $H^{s,\beta}$: only
finitely many frequency shifts occur, and the ratio of the corresponding
weights is bounded.

\subsection{The finite matrix and the positive square-mask hypothesis}
Set
\begin{equation}
 d=\left\lfloor\frac D{b-1}\right\rfloor,\qquad
 F_d=\Span\{e_j:|j|\le d\},\qquad
 A_{b,w}=(w_{bk-j})_{-d\le k,j\le d}.
 \label{eq:finite-matrix}
\end{equation}
A coefficient outside $[-D,D]$ is understood to be zero. This matrix is
exactly the restriction of $\Lop_{b,w}$ to $F_d$.

The hypothesis used in the general result is
% ed. (2026-09-29): equation* for the tagged display; the numbered environment
% gave hyperref a duplicate destination equation.2.5.
\begin{equation*}
 h\in C(\T),\quad \min_\T h>0,\quad \rho>0,\qquad
 \Lop_{b,|w|^2}h=\rho h.
 \tag{P}\label{eq:P}
\end{equation*}
It is important that $h$ be strictly positive, not merely nonnegative.
A zero of $h$ would invalidate the bounded similarity used below.
We prove \eqref{eq:P} for all masks $\sin^{2m}(\pi x)$ in \cref{sec:sine}.
For a general complex trigonometric mask, it remains an explicit hypothesis.

\begin{theorem}[Exact Sobolev disk and finite outliers]
\label{thm:general-spectrum}
Assume \eqref{eq:P}. For all real $s,\beta$, set $r_s=b^{-s}\sqrt\rho$.
Then
\begin{align}
 \sigma_{\ess}(\Lop_{b,w};H^{s,\beta})&=\Disk{r_s},
 \label{eq:general-essential}\\
 \sigma(\Lop_{b,w};H^{s,\beta})&=\Disk{r_s}\cup\sigma(A_{b,w}).
 \label{eq:general-spectrum}
\end{align}
Every generalized eigenvector at an eigenvalue $\lambda$ with $|\lambda|>r_s$
lies in $F_d$. The algebraic multiplicities and Jordan blocks at such
$\lambda$ agree with those of $A_{b,w}$.

Moreover, for every $n\ge1$ and every compact operator $K$ on $H^{s,\beta}$,
\begin{equation}
 \norm{\Lop_{b,w}^{\,n}-K}_{H^{s,\beta}\to H^{s,\beta}}\ge r_s^n.
 \label{eq:compact-obstruction}
\end{equation}
The compact operator may depend on $n$.
\end{theorem}

The theorem identifies a \emph{set-valued spectrum}, not only its outer
radius. It does not identify the point, continuous, and residual spectra
inside the Sobolev disk. On $L^2$, the proof supplies infinite-dimensional
eigenspaces for every point in the open disk. Compact perturbation preserves
the Fredholm essential spectrum, but not those individual eigenspaces; we
therefore make no analogous interior point-spectrum assertion at every $s$.

\section{The exact \texorpdfstring{$L^2$}{L2} disk}
\label{sec:L2}

\subsection{A similarity to a coisometry}
Let $M_g$ denote multiplication by a bounded function $g$. All adjoints in
this section are taken in $L^2(\T)$ with Lebesgue measure.

\begin{lemma}[Square-mask identity]
\label{lem:square-mask}
For bounded $h$,
\begin{equation}
 \Lop_{b,w}^*g=\overline w\,(g\circ T_b),\qquad
 \Lop_{b,w}M_h\Lop_{b,w}^*
       =M_{\Lop_{b,|w|^2}h},\quad T_bx=bx\pmod1.
 \label{eq:square-mask}
\end{equation}
\end{lemma}
\begin{proof}
Changing variables on each of the $b$ branches gives
$\int(\Lop_{b,w}f)\overline g=\int wf\,\overline{g\circ T_b}$,
which proves the adjoint formula. For the second identity, evaluate
$\Lop_{b,w}(h\overline w\,g\circ T_b)$ at $x$. The factor $g(x)$
does not depend on the chosen inverse branch and can be taken outside the
sum. What remains is $\Lop_{b,|w|^2}h(x)$.
\end{proof}

\begin{lemma}[Coisometry and power bound]
\label{lem:coisometry}
Under \eqref{eq:P}, the operator
\begin{equation}
 C=\rho^{-1/2}M_{h^{-1/2}}\Lop_{b,w}M_{h^{1/2}}
 \label{eq:C}
\end{equation}
is a coisometry: $CC^*=I$. Consequently,
\begin{equation}
 \norm{\Lop_{b,w}^{\,n}}_{2\to2}
 \le \left(\frac{\max h}{\min h}\right)^{1/2}\rho^{n/2}
 \qquad(n\ge0).
 \label{eq:L2-power}
\end{equation}
\end{lemma}
\begin{proof}
The square-mask identity gives
$CC^*=\rho^{-1}M_{h^{-1/2}}M_{\rho h}M_{h^{-1/2}}=I$.
Thus $C^*$ is an isometry and $\norm C=1$. The similarity
$\Lop_{b,w}=\sqrt\rho M_{h^{1/2}}CM_{h^{-1/2}}$ yields
\eqref{eq:L2-power}.
\end{proof}

\subsection{Why the whole disk is essential}
\begin{lemma}[Infinite branch defect]
\label{lem:kernel}
The kernel of $\Lop_{b,w}$ on $L^2(\T)$ is infinite-dimensional.
\end{lemma}
\begin{proof}
A nonzero trigonometric polynomial has only finitely many zeros on $\T$.
Choose a small interval $J\subset(0,1)$ on which two branch weights
$w(y_0(x))$ and $w(y_1(x))$, with $y_a(x)=(x+a)/b$, do not vanish.
For arbitrary $g\in L^2(J)$, define $f$ on these two branch intervals by
\[
 f(y_0(x))=w(y_1(x))g(x),\qquad
 f(y_1(x))=-w(y_0(x))g(x),
\]
and let $f$ vanish on the other branches and outside these intervals.
The two contributions to $\Lop_{b,w}f$ cancel. Because the weights are
bounded away from zero after shrinking $J$, this construction embeds the
infinite-dimensional space $L^2(J)$ into the kernel.
\end{proof}

\begin{proposition}[Complete essential disk on $L^2$]
\label{prop:L2-spectrum}
Under \eqref{eq:P},
\begin{equation}
 \sigma(\Lop_{b,w};L^2)=\sigma_{\ess}(\Lop_{b,w};L^2)
 =\Disk{\sqrt\rho}.
 \label{eq:L2-disk}
\end{equation}
Every point of the open disk is an eigenvalue of infinite geometric
multiplicity.
\end{proposition}
\begin{proof}
Let $E=\ker C$, which is infinite-dimensional by \cref{lem:kernel} and the
similarity. For $e\in E$, the vectors $(C^*)^je$, $j\ge0$, lie in mutually
orthogonal copies of $E$. Indeed, if $j>k$, cancellation of $k$ copies of
$CC^*=I$ reduces the inner product to one involving a positive power of $C$
applied to an element of $\ker C$, and it is zero.

For $|z|<1$, the norm-convergent series
\begin{equation}
 u_z(e)=\sum_{j=0}^{\infty}z^j(C^*)^je
 \label{eq:disk-eigenvectors}
\end{equation}
satisfies $Cu_z(e)=zu_z(e)$ and
$\norm{u_z(e)}^2=\norm e^2/(1-|z|^2)$. It is injective in $e$.
Therefore $C-zI$ has infinite-dimensional kernel and is not Fredholm.
The non-Fredholm set is closed, so it contains the closed unit disk.
Since $\norm C=1$, there is no spectrum outside that disk. Similarity and
scaling prove the claim for $\Lop_{b,w}$.
\end{proof}

No abstract classification of all weighted shifts is needed for this
argument. The branch defect and the positive square-mask eigenfunction
supply the complete disk directly.

\section{Compact conjugacy and the exclusion of hidden outliers}
\label{sec:sobolev-proof}

\subsection{The compact remainder in a change of regularity}
Let $U_{s,\beta}:H^{s,\beta}\to L^2$ be the unitary Fourier multiplier
\begin{equation}
 U_{s,\beta}e_k=q_{s,\beta}(k)e_k,\qquad
 q_{s,\beta}(k)=\br{k}^{s}[\log(e+|k|)]^\beta.
 \label{eq:unitary}
\end{equation}

\begin{lemma}[Compact Sobolev conjugacy]
\label{lem:compact-conjugacy}
For every real $s,\beta$,
\begin{equation}
 U_{s,\beta}\Lop_{b,w}U_{s,\beta}^{-1}
 =b^{-s}\Lop_{b,w}+K_{s,\beta},
 \qquad K_{s,\beta}\text{ compact on }L^2.
 \label{eq:compact-conjugacy}
\end{equation}
For every fixed positive integer $n$, it follows that
\begin{equation}
 U_{s,\beta}\Lop_{b,w}^{\,n}U_{s,\beta}^{-1}
 =b^{-ns}\Lop_{b,w}^{\,n}+K_{s,\beta,n},
 \qquad K_{s,\beta,n}\text{ compact}.
 \label{eq:compact-power}
\end{equation}
\end{lemma}
\begin{proof}
In the Fourier rule, the entry that sends input frequency $bk-\ell$ to
output frequency $k$ is multiplied under conjugacy by
$q_{s,\beta}(k)/q_{s,\beta}(bk-\ell)$. For each fixed $\ell$,
\begin{equation}
 \frac{q_{s,\beta}(k)}{q_{s,\beta}(bk-\ell)}\longrightarrow b^{-s}
 \qquad (|k|\to\infty).
 \label{eq:weight-ratio}
\end{equation}
The power part has this limit because
$\br{k}/\br{bk-\ell}\to b^{-1}$; the quotient of the two logarithms tends
to one. The difference operator for this $\ell$ is a bounded partial
frequency shift followed by a diagonal multiplier tending to zero at both
ends of $\Z$. Such a diagonal multiplier is a norm limit of finite-rank
truncations. Summing over the finitely many $\ell$ proves
\eqref{eq:compact-conjugacy}. Taking powers preserves equality modulo
compact operators and gives \eqref{eq:compact-power}.
\end{proof}

The Fredholm essential spectrum is invariant under compact perturbations.
Together with \cref{prop:L2-spectrum}, \cref{lem:compact-conjugacy}
already proves \eqref{eq:general-essential}. This step alone does
\emph{not} exclude additional isolated eigenvalues outside the disk.
The finite-frequency argument below does that work.

\subsection{Finite bands and quotient tails}
\begin{lemma}[Invariance and strict descent]
\label{lem:band}
For every $N\ge d$, $F_N$ is invariant under $\Lop_{b,w}$. If a
trigonometric polynomial has degree $p>d$, its image has degree at most
$\lfloor(p+D)/b\rfloor<p$. Therefore the induced operator on $F_N/F_d$
is nilpotent.
\end{lemma}
\begin{proof}
A nonzero output frequency satisfies $|bk-j|\le D$, so its absolute value
is at most $(|j|+D)/b$. For $N\ge d$, write
$D=(b-1)d+e$ with $0\le e<b-1$. Then
$\lfloor(N+D)/b\rfloor\le N$. If $p>d$, one has $D<(b-1)p$,
which gives strict descent. Iteration reaches $F_d$ in finitely many
steps.
\end{proof}

\begin{lemma}[Every exterior spectral point is finite-dimensional]
\label{lem:outliers}
Under \eqref{eq:P}, if $|\lambda|>r_s$, then
$\lambda\in\sigma(\Lop_{b,w};H^{s,\beta})$ if and only if
$\lambda\in\sigma(A_{b,w})$. The generalized eigenspaces at such
$\lambda$ are contained in $F_d$.
\end{lemma}
\begin{proof}
Write $\widetilde L=U_{s,\beta}\Lop_{b,w}U_{s,\beta}^{-1}$.
Let $Q_N$ be the orthogonal projection onto frequencies $|k|>N$.
By \eqref{eq:L2-power}, choose a positive integer $n$ with
\[
 b^{-ns}\norm{\Lop_{b,w}^{\,n}}_{2\to2}<|\lambda|^n.
\]
Using \eqref{eq:compact-power} and
$\norm{Q_NK_{s,\beta,n}Q_N}\to0$, choose $N\ge d$ so large that
\begin{equation}
 \norm{Q_N\widetilde L^{\,n}Q_N}<|\lambda|^n.
 \label{eq:tail-power-bound}
\end{equation}
Since $F_N$ is invariant, the quotient operator is represented on its
orthogonal complement by $T_N=Q_N\widetilde LQ_N$, and
$T_N^n=Q_N\widetilde L^{\,n}Q_N$ on that complement. Hence
$\lambda I-T_N$ is invertible: its spectral radius is smaller than
$|\lambda|$.

Relative to $F_N\oplus F_N^\perp$, the operator is block upper triangular.
The lower diagonal block is invertible after subtracting $\lambda I$.
Thus the full operator is invertible exactly when its finite block on
$F_N$ is invertible. By \cref{lem:band}, that finite block has the same
nonzero eigenvalues as the block on $F_d$; the remaining quotient is
nilpotent. Also, if $(\widetilde L-\lambda I)^jv=0$, invertibility on the
tail forces $v\in F_N$, and invertibility of the nilpotent quotient minus
$\lambda I$ forces $v\in F_d$. Diagonal conjugation by $U_{s,\beta}$
preserves each finite band, completing the proof.
\end{proof}

\begin{proof}[Completion of \cref{thm:general-spectrum}]
The essential disk follows from compact conjugacy. Every finite-matrix
eigenvalue is an eigenvalue on the full space, since $F_d$ is invariant.
\Cref{lem:outliers} excludes any other spectrum outside the disk and
identifies the generalized eigenspaces. Near an exterior eigenvalue the
quotient block remains invertible, so the finite-block resolvent also
identifies its algebraic multiplicity and Jordan structure.

Finally, the Calkin image of $\Lop_{b,w}$ has spectrum $\Disk{r_s}$.
The spectrum of its $n$th power is $\Disk{r_s^n}$ by spectral mapping.
Its norm is at least its spectral radius, and the quotient norm is no
larger than $\norm{\Lop_{b,w}^{\,n}-K}$ for any compact $K$.
This proves \eqref{eq:compact-obstruction}.
\end{proof}

\begin{remark}[Regularity-dependent, not contradictory spectra]
An analytic Hardy space and a finite-order Sobolev space are different
Banach spaces with different resolvents. The appearance of an essential
disk here neither contradicts nor weakens the analytic finite-spectrum
result in the repository. It quantifies precisely what analytic regularity
had removed.
\end{remark}

\section{All even sine powers and integer interval regularity}
\label{sec:sine}

\subsection{A positive eigenfunction exists in a finite cone}
For $m\ge1$, let
\begin{equation}
 w_m(x)=\sin^{2m}(\pi x),\qquad
 c_{m,\ell}=\frac{(-1)^\ell}{4^m}\binom{2m}{m+\ell}
 \quad(|\ell|\le m).
 \label{eq:sine-coefficients}
\end{equation}
The coefficients are zero outside this interval. Denote the matrix
\eqref{eq:finite-matrix} for $w_m$ by $A_{b,m}$.

\begin{proposition}[Automatic positivity for even sine masks]
\label{prop:sine-positive}
For every $b\ge2$ and $m\ge1$, there exist a strictly positive real
trigonometric polynomial $h$ and a number $\rho_{b,m}>0$ such that
\begin{equation}
 \Lop_{b,w_{2m}}h=\rho_{b,m}h.
 \label{eq:sine-positive}
\end{equation}
The degree of $h$ is at most $d_2=\lfloor2m/(b-1)\rfloor$, and
\begin{equation}
 \rho_{b,m}=r(A_{b,2m}),
 \label{eq:rho-finite}
\end{equation}
where $r$ denotes spectral radius. In particular, $\rho_{b,m}$ is a
positive algebraic number.
\end{proposition}
\begin{proof}
Let $V$ be the real vector space of real-valued trigonometric polynomials
of degree at most $d_2$. Consider the set
\[
 \mathcal K=\left\{f\in V:f\ge0,\ \int_0^1f=1\right\}.
\]
It is nonempty and convex. It is compact because, for $f\in\mathcal K$,
$|\widehat f(k)|\le\int f=1$, and it is closed in the finite-dimensional
space $V$. By the finite-band lemma, $\Lop_{b,w_{2m}}$ maps $V$ into $V$.
For nonzero nonnegative $f$, the integral
$\int w_{2m}f$ is strictly positive, since $w_{2m}>0$ except at the integer
point and $f$ is a nonzero continuous function. Therefore
\[
 f\longmapsto
 \frac{\Lop_{b,w_{2m}}f}{\int_0^1\Lop_{b,w_{2m}}f}
\]
is a continuous map from $\mathcal K$ to itself. The finite-dimensional
Brouwer fixed-point theorem gives an eigenfunction $h\ge0$ with mean one
and eigenvalue $\rho>0$.

It remains to prove strict positivity. A nonzero trigonometric polynomial
of degree at most $d_2$ has at most $2d_2$ distinct zeros on $\T$: multiply
its Laurent polynomial by the $d_2$th power of its variable and use the
ordinary polynomial zero bound. Choose $n$ with
$(b-1)b^{n-1}>2d_2$. If $x\in(0,1)$, all $b^n$ inverse branches in
$\Lop_{b,w_{2m}}^{\,n}h(x)$ have strictly positive product weights.
If $x=0$, at least the $(b-1)b^{n-1}$ preimages $k/b^n$ with $b\nmid k$
have strictly positive product weights; none of their first $n$ forward
positions is an integer. There are therefore more positive-weight
preimages than zeros of $h$. At least one summand is positive, at every
$x$, so $\Lop_{b,w_{2m}}^{\,n}h(x)>0$. Since this equals $\rho^nh(x)$,
we obtain $h>0$.

On the complexification of $V$, equip functions with the equivalent norm
$\norm f_h=\max_x|f(x)|/h(x)$. Positivity gives
$|\Lop_{b,w_{2m}}f|\le\Lop_{b,w_{2m}}|f|\le\rho\norm f_h h$.
Hence its finite-dimensional spectral radius is at most $\rho$, and the
eigenfunction gives equality. Its matrix has rational entries by
\eqref{eq:sine-coefficients}, proving algebraicity.
\end{proof}

\begin{corollary}[Complete even-moment spectra]
\label{cor:all-sine}
For all integers $b\ge2$, $m\ge1$, and real $s,\beta$,
\begin{align}
 \sigma_{\ess}(\Lop_{b,w_m};H^{s,\beta})
 &=\Disk{b^{-s}\sqrt{r(A_{b,2m})}},\notag\\
 \sigma(\Lop_{b,w_m};H^{s,\beta})
 &=\Disk{b^{-s}\sqrt{r(A_{b,2m})}}\cup\sigma(A_{b,m}).
 \label{eq:all-sine-spectrum}
\end{align}
An eigenvalue $\lambda\ne0$ of $A_{b,m}$ lies outside the essential disk
exactly when
\begin{equation}
 s>\log_b\!\left(\frac{\sqrt{r(A_{b,2m})}}{|\lambda|}\right).
 \label{eq:general-isolation-threshold}
\end{equation}
\end{corollary}

This gives a finite algebraic procedure for the complete spectral set at
any specified regularity. Numerical diagonalization is not part of the
proof. Exact characteristic polynomials or certified algebraic root
isolation may be used to turn the formulas into explicit numbers.

\subsection{Unmatched interval jets do not add nonzero spectrum}
Let $H^r(0,1)$ be the usual interval Sobolev space at an integer order
$r\ge0$. When $r\ge1$, its periodic subspace is characterized by the
$r$ matching conditions
\begin{equation}
 f^{(j)}(1)=f^{(j)}(0),\qquad 0\le j<r.
 \label{eq:matching-jets}
\end{equation}
It is naturally isomorphic to $H^r(\T)$, with equivalent norms. Define
$\delta_j(f)=f^{(j)}(1)-f^{(j)}(0)$.

\begin{lemma}[Exact boundary-jet action]
\label{lem:jets}
For a smooth periodic mask $w$ and $0\le j<r$,
\begin{equation}
 \delta_j(\Lop_{b,w}f)=b^{-j-1}
 \sum_{\ell=0}^j\binom j\ell w^{(j-\ell)}(0)\delta_\ell(f).
 \label{eq:jet-action}
\end{equation}
For $w=w_m$, the induced operator on the quotient by the periodic subspace
is nilpotent, of nilpotency index at most $\lceil r/(2m)\rceil$ when $r>0$.
\end{lemma}
\begin{proof}
Differentiate the branch formula $j$ times. The difference between its
values at $1$ and $0$ telescopes over the $b$ branch endpoints, leaving
$b^{-j-1}[(wf)^{(j)}(1)-(wf)^{(j)}(0)]$. Leibniz's rule and periodicity
of all derivatives of $w$ give \eqref{eq:jet-action}. The identity holds
first for smooth functions and then by Sobolev trace continuity.

For $w_m$, the derivatives of orders $0,\ldots,2m-1$ vanish at zero.
Thus the output jet of index $j$ depends only on input jets of indices
at most $j-2m$. Iterating $\lceil r/(2m)\rceil$ times annihilates all
$r$ quotient coordinates. These coordinates describe the full quotient:
the jet map is onto by elementary Hermite interpolation.
\end{proof}

\begin{theorem}[Integer-order interval spectrum]
\label{thm:interval}
For integers $r\ge0$, $b\ge2$, $m\ge1$, the complete spectrum and the
Fredholm essential spectrum of $\Lop_{b,w_m}$ on $H^r(0,1)$ are exactly
those in \eqref{eq:all-sine-spectrum} with $s=r$ and $\beta=0$.
Every generalized eigenvector at a nonzero eigenvalue has matching endpoint
jets and belongs to the periodic subspace.
\end{theorem}
\begin{proof}
For $r=0$, interval and periodic $L^2$ are the same Hilbert space up to
identification of endpoints of measure zero. For $r>0$, the periodic
subspace is invariant and has finite codimension $r$.
Choose any bounded finite-dimensional complement. The resulting block
upper triangular operator has the periodic operator on one diagonal block
and the nilpotent jet operator on the other. For $\lambda\ne0$, the
quotient block minus $\lambda I$ is invertible. The full operator is
therefore invertible exactly when the periodic block is invertible.
The point zero already belongs to the periodic essential disk. A
finite-dimensional block and the off-diagonal finite-rank block do not
change Fredholm essential spectrum. Finally, projecting a nonzero
$\lambda$-generalized eigenvector to the nilpotent quotient gives zero,
which proves the matching assertion.
\end{proof}

This argument does not justify the same statement on fractional interval
Sobolev spaces. Trace conditions change at half-integer orders, and the
critical logarithmic spaces require separate endpoint analysis. That
extension is retained as a research question rather than silently assumed.

\section{The quadratic dyadic operator: exact thresholds}
\label{sec:dyadic-spectrum}

\subsection{The two finite matrices}
For $b=2,m=1$, the terminal Fourier band is $F_1$. In the basis
$(e_{-1},e_0,e_1)$,
\begin{equation}
 A_{2,1}=\begin{pmatrix}
 -1/4&0&0\\[-1mm]
 -1/4&1/2&-1/4\\[-1mm]
 0&0&-1/4
 \end{pmatrix},\qquad
 \det(tI-A_{2,1})=(t-1/2)(t+1/4)^2.
 \label{eq:A21}
\end{equation}
The eigenspaces are generated by
\begin{equation}
 1,\qquad v_+=e_1+\tfrac13,\qquad v_-=e_{-1}+\tfrac13,
 \label{eq:primal-eigenfunctions}
\end{equation}
with respective eigenvalues $1/2,-1/4,-1/4$. Both values are semisimple
in this finite block.

The squared mask is $w_2=\sin^4(\pi x)$. In the basis
$(e_{-2},e_{-1},e_0,e_1,e_2)$ its matrix is
\begin{equation}
 A_{2,2}=\begin{pmatrix}
 1/16&0&0&0&0\\
 3/8&-1/4&1/16&0&0\\
 1/16&-1/4&3/8&-1/4&1/16\\
 0&0&1/16&-1/4&3/8\\
 0&0&0&0&1/16
 \end{pmatrix}.
 \label{eq:A22}
\end{equation}
Its characteristic polynomial is
\begin{equation}
 (t+1/4)(t-1/16)^2(t^2-t/8-1/16).
 \label{eq:A22-char}
\end{equation}
In particular,
\begin{equation}
 \rho=\frac{1+\sqrt{17}}{16},\quad
 h(x)=1+\kappa\cos(2\pi x),\quad
 \kappa=\frac{5-\sqrt{17}}4
 \label{eq:explicit-h}
\end{equation}
satisfy $\Lop_{2,w_2}h=\rho h$. Since $0<\kappa<1$, this also gives
a direct check of the strict positivity hypothesis without the cone
argument. The exact identities in this subsection are elementary matrix
computations; their role is to make the general theorem explicit, not to
claim these finite eigenvalues as newly discovered.

\begin{theorem}[Dyadic regularity transition]
\label{thm:dyadic-spectrum}
For every real $s,\beta$, the quadratic dyadic operator \eqref{eq:intro-L}
has the spectrum \eqref{eq:intro-spectrum}. Its eigenvalue $1/2$ is isolated
if and only if $s>s_0$, and $-1/4$ is isolated if and only if $s>s_1$.
When isolated, their algebraic multiplicities are respectively one and two,
and both are semisimple. The same spectral sets hold on $H^r(0,1)$ at
integer $r\ge0$.
\end{theorem}
\begin{proof}
Use \cref{thm:general-spectrum,thm:interval} and
\eqref{eq:A21}--\eqref{eq:explicit-h}. Solving
$\gamma2^{-s}<1/2$ and $\gamma2^{-s}<1/4$ gives the two strict thresholds.
At equality, the eigenvalue is on the boundary of a disk contained in the
spectrum, and is therefore not isolated.
\end{proof}

For orientation, the constants have the approximate values
\begin{equation}
 \gamma\approx0.5658569621,\qquad
 s_0\approx0.1785093184,\qquad s_1\approx1.1785093184.
 \label{eq:threshold-decimals}
\end{equation}
These decimals are explanatory approximations; the exact radicals and
logarithms, not rounded values, define the thresholds.

\begin{center}
\begin{tabular}{lll}
\toprule
Space & Essential disk radius & Isolated displayed eigenvalues\\
\midrule
$L^2$ & $\gamma$ & none\\
$H^1$ & $\gamma/2$ & $1/2$ only\\
$H^{s_1,\beta}$ & $1/4$ & $1/2$ only\\
$H^2$ & $\gamma/4$ & $1/2$ and $-1/4$\\
\bottomrule
\end{tabular}
\end{center}

The table concerns isolation, not existence: the polynomial eigenfunctions
\eqref{eq:primal-eigenfunctions} belong to all these spaces. In particular,
one full derivative is not enough to isolate the alternating eigenvalue.

\section{Dual distributions and an exact Fourier cutoff}
\label{sec:cutoff}

\subsection{Two familiar recurrences with different spectral roles}
Define real sequences $(\eta_k)_{k\ge0}$ and $(a_k)_{k\ge0}$ by
\begin{align}
 &\eta_0=1,\qquad \eta_1=-\tfrac13,\qquad
 \eta_{2k}=\eta_k,\qquad
 \eta_{2k+1}=-\tfrac12(\eta_k+\eta_{k+1}),
 \label{eq:eta-recurrence}\\
 &a_0=0,\qquad a_1=1,\qquad
 a_{2k}=-2a_k,\qquad a_{2k+1}=a_k+a_{k+1}.
 \label{eq:a-recurrence}
\end{align}
The potentially self-referential odd equations at $k=0$ are consistent
with these initial values. For subsequent indices they give an unambiguous
recursion. The sequence $\eta$ is the balanced Thue--Morse pair-correlation
sequence; $a_k=B_k(-2)$ in the standard Stern-polynomial convention
\cite{baake-coons,ulas-ulas}. Neither sequence is introduced here as a new
integer or rational sequence.

For a trigonometric polynomial $f$, define
\begin{align}
 \mu(f)&=\sum_{k\in\Z}\eta_{|k|}\widehat f(k),\notag\\
 D_+(f)&=\sum_{k\ge1}a_k\widehat f(k),\qquad
 D_-(f)=\sum_{k\le-1}a_{|k|}\widehat f(k).
 \label{eq:dual-functionals}
\end{align}
These formulas define distributions even before their sharp Sobolev
continuity thresholds are known: $|\eta_k|\le1$ by induction, and
$|a_k|\le k$ for $k\ge1$, again by induction using
\eqref{eq:a-recurrence}.

\begin{proposition}[Dual eigenrelations and biorthogonality]
\label{prop:dual-eigen}
On trigonometric polynomials,
\begin{equation}
 \mu\Lop=\tfrac12\mu,\qquad
 D_+\Lop=-\tfrac14D_+,\qquad
 D_-\Lop=-\tfrac14D_-.
 \label{eq:dual-eigen}
\end{equation}
Moreover,
\begin{equation}
 \mu(1)=1,\quad\mu(v_\pm)=0,\quad D_\pm(1)=0,
 \quad D_+(v_+)=D_-(v_-)=1,\quad D_+(v_-)=D_-(v_+)=0.
 \label{eq:biorthogonal}
\end{equation}
Consequently the algebraic operators
\begin{equation}
 P_0f=\mu(f)1,\qquad
 P_1f=D_+(f)v_++D_-(f)v_-
 \label{eq:P01}
\end{equation}
are disjoint idempotents commuting with $\Lop$ on polynomials.
\end{proposition}
\begin{proof}
The Fourier action is
\begin{equation}
 \Lop e_{2k}=\tfrac12e_k,\qquad
 \Lop e_{2k+1}=-\tfrac14(e_k+e_{k+1}).
 \label{eq:dyadic-mode-rule}
\end{equation}
For $\mu$, its evaluation against these identities is exactly
\eqref{eq:eta-recurrence}. For $D_+$ on positive indices, it is exactly
\eqref{eq:a-recurrence}; the mode at zero is annihilated, and negative
modes cannot produce a positive output frequency. The $D_-$ argument is
symmetric. Direct evaluation on $1,v_+,v_-$ gives
\eqref{eq:biorthogonal} and then the projection identities.
\end{proof}

\subsection{The finite Riesz product and its stable coefficient formula}
Write
\begin{equation}
 p_n(x)=2^nW_n(x)=\prod_{j=0}^{n-1}(1-\cos(2\pi2^jx)),
 \qquad q_n(k)=\widehat p_n(k).
 \label{eq:pn}
\end{equation}
The polynomial $p_n$ is nonnegative, has degree $2^n-1$, and has mean one.
Indeed, $p_0=1$ and the identity
$p_{n+1}(x)=(1-\cos(2\pi x))p_n(2x)$ gives
\begin{equation}
 q_{n+1}(2k)=q_n(k),\qquad
 q_{n+1}(2k+1)=-\tfrac12(q_n(k)+q_n(k+1)).
 \label{eq:qn-recurrence}
\end{equation}
In particular $q_n(0)=1$ and $q_n(-k)=q_n(k)$.

\begin{theorem}[Exact cutoff identity]
\label{thm:cutoff}
For every $n\ge0$ and $0\le k\le2^n$,
\begin{equation}
 q_n(k)=\eta_k+\frac13(-1/2)^na_k.
 \label{eq:cutoff}
\end{equation}
For $|k|\ge2^n$, $q_n(k)=0$. At $k=2^n$ both statements agree.
\end{theorem}
\begin{proof}
At $n=0$, the required indices are $k=0,1$. The right side is respectively
$1$ and $-1/3+1/3=0$, as required. Assume the formula at level $n$.
For an even index $2k\le2^{n+1}$, combine
$q_{n+1}(2k)=q_n(k)$ with $\eta_{2k}=\eta_k$ and $a_{2k}=-2a_k$;
the correction factor satisfies
$(-1/2)^{n+1}(-2)=(-1/2)^n$.
For an odd index $2k+1\le2^{n+1}$, the two indices $k,k+1$ lie in the
previous valid range. Substitution in \eqref{eq:qn-recurrence}, followed
by the two sequence recurrences, gives the formula.
The support statement follows from the degree of $p_n$. At its boundary,
$\eta_{2^n}=-1/3$ and $a_{2^n}=(-2)^n$, so the right side is zero.
\end{proof}

As a consequence, $q_n(k)\to\eta_{|k|}$ for every fixed integer $k$.
The probability measures $p_n(x)\dd x$ therefore converge weakly to a
probability measure whose integral against every trigonometric polynomial
is $\mu$. To see this without any measure-theoretic identification theorem,
use compactness of probability measures on the compact circle to obtain
subsequential limits, then use convergence on all characters and density of
trigonometric polynomials in $C(\T)$ to prove uniqueness. Thus the notation
$\mu$ can denote both the leading distribution and this measure.
No singularity assertion is needed for any proof below.

\subsection{The moment remainder is exactly a Fourier tail}
For a polynomial $f$, change of variables in the transfer formula gives
\begin{equation}
 I_n(f)=\int_0^1\Lop^nf(x)\dd x.
 \label{eq:transfer-moment}
\end{equation}
For example, induction follows from
$\int (\Lop g)h=\int wg(h\circ T_2)$.

\begin{corollary}[Exact two-term identity with a vanishing low-frequency remainder]
\label{cor:tail-formula}
On trigonometric polynomials, \eqref{eq:intro-two-term} holds with
\begin{equation}
 R_n(f)=-\sum_{|k|>2^n}
 \left(2^{-n}\eta_{|k|}+\frac{(-1/4)^n}{3}a_{|k|}\right)\widehat f(k).
 \label{eq:exact-tail}
\end{equation}
In particular, $R_n$ annihilates every polynomial of degree at most $2^n$.
Whenever $\mu,D_+,D_-$ are continuous on $H^{s,\beta}$, the identities extend
to that whole space and the tail pairing converges absolutely by
Cauchy--Schwarz.
\end{corollary}
\begin{proof}
Since $q_n$ is even, integration gives
$I_n(f)=2^{-n}\sum_kq_n(k)\widehat f(k)$. Apply
\cref{thm:cutoff} on $|k|\le2^n$, and use its support assertion outside
that interval. The coefficient identity
$2^{-n}(-1/2)^n=(-1/4)^n$ gives the result. Extension follows by density
and continuity; $I_n$ is a finite Fourier pairing on every $H^{s,\beta}$.
\end{proof}

\begin{example}[A single frequency and the nonuniformity of the cutoff]
For $f=e_{2^j}$ and $n\ge j$, the exact expression is
\[
 I_n(e_{2^j})=-\frac{2^{-n}}3+
       \frac{(-1/4)^n(-2)^j}{3}.
\]
For $n<j$, the moment is zero because its frequency lies beyond the support
of $p_n$. At $n=j$, the displayed two terms cancel exactly. This illustrates
why a fixed-frequency asymptotic is not by itself a uniform statement over
a Sobolev unit ball: frequencies can move with $n$.
\end{example}

\section{Quadratic energies and sharp continuity thresholds}
\label{sec:energies}

\subsection{Four finite-state recurrences}
Define the prefix energies and adjacent correlations
\begin{align}
 E_n&=\sum_{k=0}^{2^n-1}a_k^2,&
 C_n&=\sum_{k=0}^{2^n-1}a_ka_{k+1},\notag\\
 F_n&=\sum_{k=0}^{2^n-1}\eta_k^2,&
 G_n&=\sum_{k=0}^{2^n-1}\eta_k\eta_{k+1}.
 \label{eq:four-energies}
\end{align}
The letter $F_n$ in this section denotes a scalar energy, not the finite
Fourier space $F_N$; the indices and context distinguish them.

\begin{lemma}[Exact energy recurrences]
\label{lem:energy-recurrences}
The initial values are $(E_0,C_0)=(0,0)$ and $(F_0,G_0)=(1,-1/3)$.
For every $n\ge0$,
\begin{align}
 E_{n+1}&=6E_n+2C_n+4^n,&
 C_{n+1}&=-4E_n-4C_n-2\cdot4^n,\label{eq:EC}\\
 F_{n+1}&=\tfrac32F_n+\tfrac12G_n-\tfrac29,&
 G_{n+1}&=-F_n-G_n+\tfrac49.\label{eq:FG}
\end{align}
\end{lemma}
\begin{proof}
Put $M=2^n$. Splitting $E_{n+1}$ into even and odd indices gives
\[
 E_{n+1}=\sum_{k=0}^{M-1}\bigl[4a_k^2+(a_k+a_{k+1})^2\bigr].
\]
Since $a_0=0$ and $a_M^2=4^n$, the shifted sum
$\sum_{k=0}^{M-1}a_{k+1}^2$ equals $E_n+4^n$. This proves the first
recurrence. The adjacent products in an even--odd pair sum to
$-2(a_k+a_{k+1})^2$, proving the second.

For $\eta$, the same split uses
$\eta_{2k}=\eta_k$ and
$\eta_{2k+1}=-(\eta_k+\eta_{k+1})/2$. The shifted square sum is now
$F_n-1+1/9=F_n-8/9$, since $\eta_M^2=1/9$.
Thus
\[
 F_{n+1}=F_n+\tfrac14(2F_n+2G_n-8/9),\qquad
 G_{n+1}=-\tfrac12(2F_n+2G_n-8/9),
\]
which is \eqref{eq:FG}.
\end{proof}

\subsection{Closed forms and dyadic annuli}
Let $u=\sqrt{17}$, $\Lambda_\pm=1\pm u$, and
$\lambda_\pm=(1\pm u)/4$.

\begin{proposition}[Closed quadratic-energy formulas]
\label{prop:closed-energies}
For every $n\ge0$,
\begin{align}
 E_n&=\frac{5+u}{4u}\Lambda_+^n
      +\frac{u-5}{4u}\Lambda_-^n-\frac12\,4^n,
 \label{eq:Eclosed}\\
 F_n&=\frac{19+5u}{18u}\lambda_+^n
      +\frac{5u-19}{18u}\lambda_-^n+\frac49.
 \label{eq:Fclosed}
\end{align}
Consequently, as $j\to\infty$,
\begin{align}
 \sum_{k=2^j}^{2^{j+1}-1}a_k^2
 &\sim c_a\Lambda_+^j,&c_a&=\frac{5+u}{4},\label{eq:a-annulus}\\
 \sum_{k=2^j}^{2^{j+1}-1}\eta_k^2
 &\sim c_\eta\lambda_+^j,&c_\eta&=\frac{19+5u}{18u}\frac{u-3}{4}>0.
 \label{eq:eta-annulus}
\end{align}
\end{proposition}
\begin{proof}
The homogeneous matrix in \eqref{eq:EC} is
$\bigl(\begin{smallmatrix}6&2\\-4&-4\end{smallmatrix}\bigr)$,
with eigenvalues $1\pm\sqrt{17}$. A particular solution has
$E_n=-4^n/2$. Solving for the two homogeneous coefficients from $E_0=0$
and $E_1=1$ gives \eqref{eq:Eclosed}.
The homogeneous matrix in \eqref{eq:FG} has eigenvalues
$(1\pm\sqrt{17})/4$, and a stationary solution has $F_n=4/9$.
Using $F_0=1$, $F_1=10/9$ gives \eqref{eq:Fclosed}. These computations
can alternatively be checked by substituting into the recurrences and
initial values, which also proves them by induction.

For $E_n$, the growth rate $\Lambda_+=1+u$ is strictly larger than
$4$ and $|\Lambda_-|$. For $F_n$, $\lambda_+>1$ and
$|\lambda_-|<1$. Taking differences between consecutive prefix energies
therefore gives the annulus equivalents and the stated positive constants.
\end{proof}

Notice the exact relations
\begin{equation}
 \Lambda_+=2^{2s_1},\qquad \lambda_+=2^{2s_0},\qquad
 \Lambda_+=4\lambda_+.
 \label{eq:energy-exponents}
\end{equation}
Thus the same two thresholds that govern spectral isolation also govern
the growth of the dual coefficient energies. This agreement is derived,
not assumed from a general regularity principle.

\subsection{The Hilbert dual criterion}
\begin{lemma}[Coefficient criterion for a continuous functional]
\label{lem:dual-criterion}
Let $T(f)=\sum_kt_k\widehat f(k)$ on polynomials. Then $T$ extends
continuously to $H^{s,\beta}$ if and only if
\begin{equation}
 \sum_{k\in\Z}|t_k|^2\br{k}^{-2s}
       [\log(e+|k|)]^{-2\beta}<\infty.
 \label{eq:dual-criterion}
\end{equation}
When finite, the sum is exactly $\norm T_{(H^{s,\beta})^*}^2$.
\end{lemma}
\begin{proof}
Cauchy--Schwarz proves sufficiency and the upper norm bound.
For necessity and equality, use finite-support coefficients proportional
to $\overline{t_k}q_{s,\beta}(k)^{-2}$ and normalize their Sobolev norm.
The resulting functional values are the square roots of the finite partial
sums in \eqref{eq:dual-criterion}. Let the support increase.
\end{proof}

\begin{theorem}[Sharp leading and subleading regularity thresholds]
\label{thm:dual-thresholds}
The functional $\mu$ extends continuously to $H^{s,\beta}$ exactly when
\begin{equation}
 s>s_0\quad\text{or}\quad s=s_0\ \text{and}\ \beta>1/2.
 \label{eq:mu-threshold}
\end{equation}
Each of $D_+$ and $D_-$, and also their sum, extends continuously exactly
when
\begin{equation}
 s>s_1\quad\text{or}\quad s=s_1\ \text{and}\ \beta>1/2.
 \label{eq:D-threshold}
\end{equation}
In their stated ranges, $P_0$ and $P_1$ are bounded commuting projections.
When the corresponding eigenvalue is isolated, these are its Riesz
projections.
\end{theorem}
\begin{proof}
On the dyadic block $2^j\le|k|<2^{j+1}$, the reciprocal squared Sobolev
weight is comparable, with constants independent of $j\ge1$, to
$2^{-2sj}(1+j)^{-2\beta}$. The annulus formulas show that the positive
frequency contributions in \eqref{eq:dual-criterion} are comparable to
\begin{equation}
 2^{2(s_0-s)j}(1+j)^{-2\beta}\quad\text{for }\mu,
 \qquad
 2^{2(s_1-s)j}(1+j)^{-2\beta}\quad\text{for }D_+.
 \label{eq:block-dual-tests}
\end{equation}
Negative frequencies give an identical test. The geometric series
converges above the stated exponent, diverges below it, and at equality
reduces to the power-series test $\sum j^{-2\beta}<\infty$.
For $D_++D_-$, the supports are disjoint, so there is no cancellation in
the squared dual norm. This proves necessity as well as sufficiency.

Bounded projection identities follow from \cref{prop:dual-eigen} and
density. When an eigenvalue is outside the essential disk, its generalized
eigenspace is exactly its semisimple finite Fourier eigenspace by
\cref{thm:general-spectrum}. On every polynomial, sufficiently many
iterations descend into $F_1$, where \eqref{eq:biorthogonal} gives the
finite spectral projection. Thus the displayed projection agrees on a
dense subspace with the Riesz projection, and hence everywhere.
\end{proof}

\begin{remark}[Continuous embedded projections]
At $s=s_1$, $\beta>1/2$, the finite-rank operator $P_1$ is bounded and
commutes with $\Lop$, but $-1/4$ is not isolated. It is therefore not a
Riesz projection defined by a contour separating that point from the
rest of the spectrum. A bounded projection onto a finite-dimensional
invariant summand is compatible with other spectrum accumulating at the
same eigenvalue. The corresponding observation applies to $P_0$ at
$s=s_0$, $\beta>1/2$.
\end{remark}

\section{Sharp scalar remainders, including the critical endpoint}
\label{sec:sharp-remainders}

\subsection{Weighted tails of the two coefficient sequences}
For $z=(z_k)_{k\ge1}$, put
\begin{equation}
 \mathcal T_z(N;s,\beta)^2=
 \sum_{k>N}|z_k|^2\br{k}^{-2s}[\log(e+k)]^{-2\beta}.
 \label{eq:weighted-tail}
\end{equation}
We use $X_n\asymp Y_n$ to mean that there exist positive constants,
independent of sufficiently large $n$, bounding $X_n/Y_n$ above and below.
Constants may depend on the fixed displayed parameters.

\begin{lemma}[Two-sided tail estimates]
\label{lem:tails}
Let $\sigma_a=s_1$ and $\sigma_\eta=s_0$. For $z=a$ or $z=\eta$,
if $s>\sigma_z$, then
\begin{equation}
 \mathcal T_z(2^n;s,\beta)
 \asymp_{s,\beta}2^{(\sigma_z-s)n}(1+n)^{-\beta}.
 \label{eq:tail-supercritical}
\end{equation}
If $s=\sigma_z$ and $\beta>1/2$, then
\begin{equation}
 \mathcal T_z(2^n;\sigma_z,\beta)
 \asymp_\beta(1+n)^{1/2-\beta}.
 \label{eq:tail-critical}
\end{equation}
The analogous norm restricted to one block $2^j\le k<2^{j+1}$ is
comparable to $2^{(\sigma_z-s)j}(1+j)^{-\beta}$.
\end{lemma}
\begin{proof}
Apply the annulus estimates in \cref{prop:closed-energies} and the
comparability of the weights on each dyadic block. The squared tail is
comparable to
\[
 \sum_{j\ge n}2^{2(\sigma_z-s)j}(1+j)^{-2\beta}.
\]
Deleting the single boundary term $k=2^n$ does not change the comparison.
For $a$, its square is $4^n=o(2^{2s_1n})$; for $\eta$, it is $1/9$
while the annular energy grows like $2^{2s_0n}$. When $s>\sigma_z$,
a decaying geometric factor makes the sum comparable to its initial
term, even for negative $\beta$. At $s=\sigma_z$, the integral test
makes the sum comparable to $(1+n)^{1-2\beta}$. Taking square roots
proves all the assertions.
\end{proof}

\subsection{The exact rate above critical regularity}
\begin{theorem}[Optimal supercritical scalar remainder]
\label{thm:supercritical-remainder}
For every $s>s_1$ and every real $\beta$, the expansion
\eqref{eq:intro-two-term} defines continuous functionals on $H^{s,\beta}$,
and
\begin{equation}
 \norm{R_n}_{(H^{s,\beta})^*}
 \asymp_{s,\beta}(\gamma2^{-s})^n(1+n)^{-\beta}.
 \label{eq:sharp-supercritical}
\end{equation}
The upper estimate is uniform for $f$ in the Sobolev unit ball, and its
exponential rate and displayed power of $n$ cannot be improved uniformly.
In particular, $\norm{R_n}=o(4^{-n})$ in this range.
\end{theorem}
\begin{proof}
Continuity follows from \cref{thm:dual-thresholds}. By the exact tail formula
and the dual norm criterion, $\norm{R_n}$ is the weighted $\ell^2$ norm of
its coefficient sequence. The triangle inequality and \cref{lem:tails}
give an upper bound by a fixed constant times
\begin{align*}
 &2^{-n}2^{(s_0-s)n}(1+n)^{-\beta}
  +4^{-n}2^{(s_1-s)n}(1+n)^{-\beta}\\
 &\hspace{4em}=2(\gamma2^{-s})^n(1+n)^{-\beta},
\end{align*}
where the harmless factor $1/3$ in the second coefficient is absorbed into
the constant. The equality of the two exponential rates uses
$2^{s_0}/2=2^{s_1}/4=\gamma$.

For the lower bound, cancellation between the two coefficient sequences
must be excluded; separate upper estimates would not suffice. Fix a large
integer $L\ge1$ and restrict to the block
$2^{n+L}\le k<2^{n+L+1}$. Let $A_{n,L}$ and $B_{n,L}$ be the weighted
$\ell^2$ norms of, respectively, the Stern term and the correlation term
on this block. The one-block estimates give
\begin{align*}
 A_{n,L}&\asymp 4^{-n}2^{(s_1-s)(n+L)}(1+n+L)^{-\beta},\\
 B_{n,L}&\asymp 2^{-n}2^{(s_0-s)(n+L)}(1+n+L)^{-\beta}.
\end{align*}
The implicit comparison constants can be chosen independently of $n,L$
when $n+L$ is sufficiently large. Hence
$A_{n,L}/B_{n,L}\ge c_{s,\beta}2^L$ for some positive constant.
Choose and fix $L$ so that this lower bound exceeds two. The reverse
triangle inequality then makes the norm of the sum at least
$A_{n,L}-B_{n,L}\ge A_{n,L}/2$. With this fixed $L$, that lower bound
is a positive constant times the right side of
\eqref{eq:sharp-supercritical}. The dual norm criterion identifies this
coefficient norm with the exact supremum over the unit ball, proving
optimality. Finally, $\gamma2^{-s}<1/4$ when $s>s_1$.
\end{proof}

\subsection{A sharp logarithmic rescue at the spectral boundary}
\begin{theorem}[Critical two-term scalar asymptotic]
\label{thm:critical-remainder}
At $s=s_1$ and $\beta>1/2$, the two-term expansion
\eqref{eq:intro-two-term} is valid on all of $H^{s_1,\beta}$ and
\begin{equation}
 \norm{R_n}_{(H^{s_1,\beta})^*}
 \asymp_\beta4^{-n}(1+n)^{1/2-\beta}.
 \label{eq:sharp-critical}
\end{equation}
Equivalently,
\begin{equation}
 2^n I_n(f)=\mu(f)+\frac{(-1/2)^n}{3}(D_++D_-)(f)
 +O_\beta\!\left(2^{-n}(1+n)^{1/2-\beta}\norm f_{s_1,\beta}\right).
 \label{eq:critical-normalized}
\end{equation}
The error is uniformly $o(2^{-n})$ after this normalization. The power
$(1+n)^{1/2-\beta}$ is sharp in the dual norm.
\end{theorem}
\begin{proof}
The coefficient functionals are bounded by \cref{thm:dual-thresholds}.
The weighted norm of the Stern part of the exact tail is comparable to
$4^{-n}(1+n)^{1/2-\beta}$ by the critical tail estimate. The correlation
part, which is one derivative above its own critical index, has norm
comparable to
\[
 2^{-n}2^{(s_0-s_1)n}(1+n)^{-\beta}
 =4^{-n}(1+n)^{-\beta}.
\]
The second quantity is smaller than the first by a factor comparable to
$(1+n)^{-1/2}$. Thus both the triangle and reverse triangle inequalities
give \eqref{eq:sharp-critical} for sufficiently large $n$. The two
frequency halves are disjoint and symmetric, contributing only a fixed
factor of $\sqrt2$. Multiplication by $2^n$ gives
\eqref{eq:critical-normalized}; since $\beta>1/2$, its extra power tends
to zero.
\end{proof}

\begin{corollary}[Necessity of the logarithmic threshold]
\label{cor:no-endpoint-without-log}
On $H^{s_1,\beta}$ with $\beta\le1/2$, there do not exist bounded linear
functionals $M,S$ and continuous remainders $\varepsilon_n$ satisfying
\[
 I_n=2^{-n}M+(-1/4)^nS+\varepsilon_n,\qquad
 \norm{\varepsilon_n}=o(4^{-n}).
\]
The same impossibility holds on $H^{s,\beta}$ for every $s<s_1$.
\end{corollary}
\begin{proof}
For a fixed polynomial, the remainder vanishes once $2^n$ exceeds its
degree. Thus the leading and subleading coefficients are uniquely fixed
by limits: the leading one is $\lim_n2^nI_n(f)=\mu(f)$, and after its
subtraction the rescaled subleading coefficient is
$(D_++D_-)(f)/3$. Any continuous coefficient functional agreeing with
those polynomial limits would extend $D_++D_-$. The necessity part of
\cref{thm:dual-thresholds} says that such an extension does not exist in
the stated range.
\end{proof}

This is an obstruction to a \emph{uniform continuous-coefficient theorem}
on the whole Hilbert space. It does not say that a particular smoother
amplitude, or a polynomial, fails to have a two-term expansion. Nor does
it rule out different formulations using unbounded coefficient maps on
proper domains.

\section{Scalar asymptotics versus operator asymptotics}
\label{sec:operator-scalar}

\subsection{The supercritical operator expansion}
\begin{proposition}[Operator remainder and its exact exponential scale]
\label{prop:operator-expansion}
For $s>s_1$ and any real $\beta$, put
$Q=\Lop(I-P_0-P_1)$ on $H^{s,\beta}$. For $n\ge1$,
\begin{equation}
 \Lop^n=2^{-n}P_0+(-1/4)^nP_1+Q^n,
 \qquad \sigma(Q)=\Disk{\gamma2^{-s}}.
 \label{eq:operator-expansion}
\end{equation}
Consequently,
\begin{equation}
 \lim_{n\to\infty}\norm{Q^n}^{1/n}=\gamma2^{-s},\qquad
 \norm{Q^n}\le C_\varepsilon(\gamma2^{-s}+\varepsilon)^n
 \quad(\varepsilon>0).
 \label{eq:operator-rate}
\end{equation}
Also $\norm{Q^n}\ge(\gamma2^{-s})^n$ for every $n\ge1$.
\end{proposition}
\begin{proof}
The bounded disjoint projections commute with $\Lop$ and have the
prescribed eigenvalues on their ranges. They are the Riesz projections
for the two isolated spectral values. Removing their summands leaves
exactly the essential disk; adding zero on their finite-dimensional
ranges does not enlarge it. This proves \eqref{eq:operator-expansion}.
The spectral radius formula proves \eqref{eq:operator-rate}. The lower
bound follows either from that spectral set and spectral mapping at
each $n$, or from \eqref{eq:compact-obstruction} after subtracting the
two finite-rank terms.
\end{proof}

The general spectral argument yields an arbitrarily small exponential
slack $\varepsilon$ in an upper operator-norm bound. The scalar moment
calculation in \cref{thm:supercritical-remainder} is stronger for its
specific observable: it removes that slack and supplies the exact
logarithmic power. One must not promote that scalar estimate to the full
operator norm without a separate proof.

\subsection{At critical regularity the two notions separate}
\begin{theorem}[A scalar little-oh law without an operator little-oh law]
\label{thm:critical-separation}
On $H^{s_1,\beta}$ with $\beta>1/2$, the operators $P_0,P_1$ remain
bounded and the algebraic power identity in \eqref{eq:operator-expansion}
remains valid. Nevertheless,
\begin{equation}
 \norm{\Lop^n-2^{-n}P_0-(-1/4)^nP_1}\ge4^{-n}
 \qquad(n\ge1).
 \label{eq:critical-operator-obstruction}
\end{equation}
For the integration functional $\ell_0(f)=\int_0^1f(x)\dd x$,
\begin{equation}
 \norm{\ell_0\bigl(\Lop^n-2^{-n}P_0-(-1/4)^nP_1\bigr)}
 \asymp_\beta4^{-n}(1+n)^{1/2-\beta}=o(4^{-n}).
 \label{eq:critical-scalar-separation}
\end{equation}
\end{theorem}
\begin{proof}
The projection identities were established algebraically and extend by
boundedness, so the power decomposition remains true without an isolation
assumption. Both subtracted terms have finite rank. The essential disk
has radius $1/4$ at $s=s_1$, and the compact obstruction in
\cref{thm:general-spectrum} proves
\eqref{eq:critical-operator-obstruction}. On the other hand,
$\ell_0(1)=1$ and $\ell_0(v_+)=\ell_0(v_-)=1/3$.
Therefore the functional on the left of
\eqref{eq:critical-scalar-separation} is exactly $R_n$.
Apply \cref{thm:critical-remainder}.
\end{proof}

The source of the improvement is visible in \eqref{eq:exact-tail}:
integration after $n$ transfer steps annihilates the remainder on every
frequency up to $2^n$. The operator itself has no comparable finite-rank
annihilation of its full essential component. The critical scalar law is
therefore not an exception to the spectral obstruction, but a different
kind of asymptotic statement.

\subsection{How the result relates to Rvachev Fourier shells}
The repository uses the entire product
\begin{equation}
 \Phi(t)=\prod_{j=0}^{\infty}\sinc\!\left(\frac{\pi t}{2^j}\right),
 \qquad \sinc z=\frac{\sin z}{z},\quad\sinc0=1.
 \label{eq:Phi}
\end{equation}
This is its Fourier-product normalization for the Rvachev up-function
\cite{repo-fourier,arias}. Its finite-scale identity, obtained simply
by separating the first $n$ factors, is
\begin{equation}
 \Phi(2^nx)=\Phi(x)\prod_{j=1}^n
       \frac{\sin(\pi2^jx)}{\pi2^jx}.
 \label{eq:Phi-scale}
\end{equation}
For real $x\ne0$, squaring and clearing the denominators gives
\begin{equation}
 (\pi x)^{2n}2^{n(n+1)}|\Phi(2^nx)|^2
       =|\Phi(x)|^2W_n(2x).
 \label{eq:Phi-Wn}
\end{equation}
This exhibits directly the dyadic sine-product moments underlying shell
analysis. The repository's precise RMS normalization packages the remaining
factors into an analytic amplitude; that analytic application and its
certified constants are credited to the antecedent manuscript, not
recomputed or claimed as new here.

The new information is what remains true after weakening regularity.
For periodic amplitudes, \cref{thm:supercritical-remainder,thm:critical-remainder}
give sharp uniform two-term laws. For nonperiodic amplitudes on an interval,
\cref{thm:interval} gives the exact integer-order spectral picture: two
Sobolev derivatives suffice to isolate both displayed dyadic eigenvalues,
whereas one does not. A nonperiodic analytic amplitude must not be declared
to lie in every periodic Sobolev space; mismatched endpoint values or jets
can prevent that. The fractional periodic endpoint theorem consequently
requires its periodic hypothesis when applied to an amplitude.

\section{Verification, implementation, and formalization boundaries}
\label{sec:verification}

\subsection{What the accompanying computation checks}
The delivered program \texttt{verify.py} uses only Python's standard
library. It constructs $\eta_k$ and $a_k$ exactly, forms the finite Riesz
products by independent Laurent-polynomial multiplication, and compares
their coefficients with \eqref{eq:cutoff}. It checks the energy recurrences
against direct finite sums and the closed forms in the exact field
$\Q(\sqrt{17})$, represented as pairs of rational numbers. It also checks
the dual eigenrelations, finite-band invariance and strict descent, the
two displayed characteristic polynomials, and the explicit positive
square-mask eigenfunction.

The executed default run, through dyadic level $12$, returned
\texttt{PASS}. Its recorded checks are:
\begin{center}
\begin{tabular}{lr}
\toprule
Checked identity or property & Number of checks\\
\midrule
Exact cutoff coefficients, including both signs & 16,395\\
Dual eigenfunctional identities & 24,579\\
Finite-band invariance and descent cases & 89,014\\
Direct energy recurrences & 52\\
Closed energy formulas in $\Q(\sqrt{17})$ & 26\\
Positive eigenfunction coordinates & 5\\
Characteristic polynomial identities & 2\\
\bottomrule
\end{tabular}
\end{center}
The frequency checks reach $|k|=4096$; the band tests cover
$b=2,\ldots,8$ and $m=1,\ldots,6$. These finite ranges are computational
scope statements, not substitutes for quantifiers in the theorems.

% ed. (2026-09-29): the default output is now rerun/verification.json, so the
% command below no longer overwrites the recorded verification.json.
Run, from the extracted directory,
\begin{verbatim}
python verify.py
\end{verbatim}
which writes \texttt{rerun/verification.json}; the recorded run is
\texttt{verification.json}.
The output file includes the exact energy data, arithmetic type, Python
version, and check counts. Explicit exceptions enforce the checks even
if Python assertions are disabled. No numerical eigensolver, floating-point
comparison, or unvalidated quadrature enters a passed mathematical check.
The approximate decimals in \eqref{eq:threshold-decimals} are not used by
the verification program.

\subsection{What is proved by analysis, not by sampling}
The infinite-dimensional spectrum is proved through the coisometry
similarity, construction of infinitely many defect eigenvectors, compact
conjugacy, and the quotient-tail argument. The endpoint norm equivalents
are proved through exact energy recurrences, dyadic comparisons, and the
Hilbert dual criterion. No finite list of Fourier coefficients could
establish those conclusions by itself.

The external functional-analytic inputs are standard: invariance of the
Fredholm essential spectrum under compact perturbations, its description
through the Calkin algebra, spectral mapping and the spectral radius
formula, finite-dimensional Brouwer fixed point, and the integer-order
Sobolev trace theorem. The special operators and estimates needed here
are proved in the text. The dependency structure is:
\begin{center}\small
\begin{tabularx}{\textwidth}{@{}lX@{}}
\toprule
Conclusion & Main ingredients\\
\midrule
Exact $L^2$ disk & Square-mask identity, strict positivity, infinite branch defect\\
Sobolev disk & Fourier weight ratio and compactness of vanishing diagonals\\
Complete outlier classification & Finite-band descent and invertible quotient tails\\
Even sine family & Finite cone, zero count, and Brouwer fixed point\\
Interval theorem & Exact boundary jets and finite-codimensional nilpotent quotient\\
Critical scalar law & Cutoff identity, exact energies, and weighted dual norms\\
\bottomrule
\end{tabularx}
\end{center}

\subsection{A realistic route to formalization}
A first formalization layer should encode the two sequences over $\Q$ and
$\Z$, prove their recurrences, and formalize the finite coefficient identity.
These statements require only finite sums, parity decompositions, and
induction. The quadratic-energy formulas can be formalized either in a
quadratic field or by verifying their second-order recurrences without
introducing radicals until the final corollary.

A second layer should prove the dyadic weighted $\ell^2$ criterion and its
geometric and critical power-series tests. That layer already yields the
sharp continuity thresholds and scalar remainder results independently of
Calkin-algebra infrastructure. The operator-spectrum layer then requires
bounded Fourier multipliers, compact operators, quotient spectra, and
Fredholm theory. It is logically separable and materially more demanding.

No new Lean code or proof objects are supplied, and no Lean build was run.
The Python interpreter, its arithmetic implementation, and the supplied
program have not been formally verified as part of this package. The
mathematical proofs and the finite checks are complementary evidence,
not interchangeable forms of certification.

\section{Further research questions}
\label{sec:questions}
The following questions are not conclusions of this paper. They separate
natural next steps from claims already established above.

\begin{question}[Fractional interval and trace-critical spaces]
Determine the exact spectrum on $H^s(0,1)$ and on logarithmically refined
interval spaces when $s$ is nonintegral. What occurs at half-integer trace
thresholds, and what are the sharp two-term moment estimates for amplitudes
whose first unmatched endpoint jet is prescribed? The finite integer-order
jet quotient suggests the correct decomposition but does not settle the
fractional norm comparison.
\end{question}

\begin{question}[Dual distributions for higher even moments]
For general $b,m$, each nonzero eigenvector of $A_{b,m}$ should lead to
coefficient functionals governed by finite recurrences along residue
classes modulo $b$. Develop an effective construction, and determine their
exact Sobolev and logarithmic thresholds. In the presence of Jordan blocks,
identify the additional powers of $n$ in moment expansions and the endpoint
logarithmic exponents they force.
\end{question}

\begin{question}[Which masks admit a positive square-mask eigenfunction?]
Characterize the nonzero trigonometric masks satisfying \eqref{eq:P}.
For even sine powers, one-point vanishing and a finite zero count suffice.
For masks whose zeros contain periodic inverse-branch configurations,
strict positivity may fail. Determine the resulting essential spectral set
and whether a decomposition into invariant components replaces the single
coisometry similarity.
\end{question}

\begin{question}[Beyond polynomial masks]
For analytic or sufficiently smooth nonpolynomial masks, Fourier conjugacy
still suggests a compact leading comparison when the coefficient tails are
controlled. Can one recover an exact Sobolev essential disk and classify
all exterior eigenvalues through a convergent infinite Fourier reduction?
The finite nilpotent quotient used here no longer exists, so its role must
be replaced rather than assumed.
\end{question}

\begin{question}[Nonlinear expanding maps]
Replace $T_bx=bx$ by a nonlinear expanding covering of the circle. The
Sobolev multiplier ratio then depends on position through the derivative.
Find sharp hypotheses under which pressure or a positive square-mask
transfer equation determines the entire essential spectral set, not merely
an upper bound on its radius. The present proof uses constant expansion
in an essential way.
\end{question}

\begin{question}[Which scalar observables improve at an embedded boundary?]
Classify continuous functionals $\ell$ for which
$\ell(\Lop^n-2^{-n}P_0-(-1/4)^nP_1)=o(4^{-n})$ at critical regularity.
Integration has an exact low-frequency cutoff. Is a quantitative variant
of this cutoff necessary, or can cancellation in other observables give
a genuinely different mechanism?
\end{question}

\begin{question}[Exact constants and dyadic oscillations in the sharp norms]
The present estimates determine the exact exponential rate and logarithmic
power up to positive constants. Determine whether the corresponding
normalized dual norms converge, and when log-periodic oscillations occur
for non-dyadic cutoffs. Such a refinement requires weighted distribution
information within each annulus, not only its total energy.
\end{question}

\begin{question}[Quantitative spectral projections near the thresholds]
As $s\downarrow s_0$ or $s\downarrow s_1$, the relevant dual norm diverges.
The block tests predict the order of this divergence. Determine exact
leading constants, bounds uniform in a joint limit of $s,\beta,n$, and
certified resolvent estimates suitable for stable numerical computation.
This would turn the sharp qualitative transition into an effective
conditioning theory.
\end{question}

\begin{question}[Finite-smoothness spaces outside the Hilbert scale]
Answer the original repository question on H\"older and $C^r$ spaces.
The Hilbert-space coisometry and exact weighted $\ell^2$ dual norms are
specific to the present proof. Their replacements must account for the
relevant Banach norms and cannot be obtained by merely replacing $s$ with
an integer differentiability index in the final formula.
\end{question}
% ed. (2026-09-29): context from the repository's rewrite of the synthesis.
\begin{ednote}
The repository's separately supplied rewrite
\path{rvachev_up_fourier_decay/Rvachev_Up_Fourier_Decay-2/} imports the
Ruelle--Perron--Frobenius theorem on $C^\alpha[0,1]$, $0<\alpha\le1$, for
this operator: there the leading eigenvalue $1/2$ is simple with a spectral
gap for every such $\alpha$. This does not conflict with
\cref{thm:dyadic-spectrum}, where $1/2$ is not isolated for $s\le s_0$: the
two statements concern different norms, $L^2$-based Sobolev norms here and
sup-norm H\"older norms there. The question above remains open.
\end{ednote}

\begin{question}[A proof-producing computational interface]
Construct a verified procedure that takes an integer dilation and a rational
trigonometric mask, produces the finite matrices and their invariant bands,
and checks a supplied positive square-mask eigenfunction certificate.
Combining this with a formal proof of \cref{thm:general-spectrum} would
turn finite algebraic data into machine-checked infinite-dimensional
spectral conclusions for a broad family of explicit examples.
\end{question}

\section{Conclusion}
The analytic finite spectrum is only one regularity regime of the
Rvachev--Thue--Morse transfer operator. On the periodic Sobolev scale,
its replacement is exact and computable: a regularity-dependent essential
disk together with a finite set of algebraic outliers. The squared sine
mask determines the disk radius, while the unsquared mask determines the
finite outliers. These are different pieces of spectral data and require
different arguments.

In the dyadic quadratic case, the leading Thue--Morse measure and the
subleading Stern distribution have sharp thresholds $s_0$ and $s_1=s_0+1$.
Their exact energy growth explains both spectral isolation and the
continuity of asymptotic coefficient functionals. The Fourier-cutoff
identity goes further: it identifies the moment remainder as a literal
high-frequency tail. This permits a uniform scalar asymptotic at a
critical essential boundary, provided and only provided that the
subleading functional is made continuous by the logarithmic exponent
$\beta>1/2$. The simultaneous operator-norm obstruction shows why that
conclusion could not have been inferred from a spectral gap.

\appendix
\section{Repository provenance and claim boundaries}
\label{sec:provenance}
The directly inspected repository materials were the Fabius-function
README, the Fourier-decay group README, the spectral-collapse manuscript
README and inventory, and the spectral manuscript source at the immutable
commit below. The source search located the explicit question
\emph{Spectra at finite smoothness}. Inspection was targeted to the relevant
research branch, not an exhaustive audit of the whole repository.

\begin{center}\small
\begin{tabularx}{\textwidth}{@{}lX@{}}
\toprule
Item & Identifier or scope\\
\midrule
Repository & \texttt{VladimirReshetnikov/ProveIt}\\
Immutable comparison commit &
\texttt{b2aed2d23125006c48d63fc97aef0d64b67becb1}\\
Spectral source blob in inspected inventory &
\texttt{1289ae8ca02383bd2a44ad8f9edb429941d004ff}\\
Spectral README blob &
\texttt{7d32919e68332fd8ca78211811b2f005ab97d175}\\
Fourier group README blob &
\texttt{30d7898b50f2bf0ea199cf24f5b543d8b98fe4fe}\\
Source question & Spectra at finite smoothness\\
Answered here & Periodic $H^{s,\beta}$, all real $s,\beta$; interval $H^r$, integer $r\ge0$\\
Not answered here & H\"older/$C^r$ spectrum; general fractional interval spaces; full interior point-spectrum decomposition\\
\bottomrule
\end{tabularx}
\end{center}
The spectral source path, relative to the repository root, is
\begin{quote}\small\ttfamily\raggedright
Analysis/FabiusFunction/docs/\allowbreak
semi-formalized-research-frontiers/drafts/\allowbreak
rvachev\_up\_fourier\_decay/\allowbreak
Spectral\_Collapse\_Alternating\_RMS/article.tex
\end{quote}
The observed commit timestamp is 30 September 2026 at 03:29:35 UTC,
which is 29 September 2026 at 20:29:35 Pacific daylight time. This accounts
for the different UTC and local calendar dates without treating either
as a future source.

The earlier manuscript's analytic determinant and RMS constant enclosures
are antecedents. They were not independently recertified here. The present
proofs do not rely on their correctness for the new Sobolev or endpoint
conclusions. Global priority and peer-review acceptance are not established.
The supplied text is intended to be mathematically checkable and is
accompanied by an explicit statement of the analytic inputs and finite
verification scope.

\section{Elementary implementation details}
\label{app:implementation}
The exact field operations for $a+b\sqrt{17}$ use
\[
 (a,b)+(c,d)=(a+c,b+d),\qquad
 (a,b)(c,d)=(ac+17bd,ad+bc).
\]
All four coordinates are rational numbers. Thus checking
\eqref{eq:Eclosed} and \eqref{eq:Fclosed} requires no approximate radical.
The two characteristic polynomials are checked by the
Faddeev--LeVerrier recurrence over $\Q$, followed by a zero-matrix
Cayley--Hamilton check. The finite Riesz products are constructed by
multiplying their Laurent coefficient dictionaries, independently of the
stable-cutoff formula.

The verification log and JSON output are part of the delivered archive.
The defaults are intentionally finite and reproducible. Increasing the
level parameter increases the amount of exact arithmetic but does not
alter the logical role of computation: the general recurrences and
functional-analytic theorems are established by proofs in the article.

\begin{thebibliography}{99}

\bibitem{repo-spectral}
Vladimir Reshetnikov, \emph{ProveIt repository},
\emph{Spectral Collapse and Alternating RMS Asymptotics for the Rvachev
Up-Function}, research manuscript dated 28 September 2026, with README.
Inspected at commit
\texttt{b2aed2d23125006c48d63fc97aef0d64b67becb1}.
\href{https://github.com/VladimirReshetnikov/ProveIt/tree/b2aed2d23125006c48d63fc97aef0d64b67becb1/Analysis/FabiusFunction/docs/semi-formalized-research-frontiers/drafts/rvachev_up_fourier_decay/Spectral_Collapse_Alternating_RMS}{Immutable repository directory}.
The source question and blob identifiers are recorded in
\cref{sec:provenance}.

\bibitem{repo-fourier}
Vladimir Reshetnikov, \emph{ProveIt repository}, Fabius-function Fourier-decay
group and canonical synthesis, with group README.
\href{https://github.com/VladimirReshetnikov/ProveIt/tree/b2aed2d23125006c48d63fc97aef0d64b67becb1/Analysis/FabiusFunction/docs/semi-formalized-research-frontiers/drafts/rvachev_up_fourier_decay}{Immutable Fourier-decay directory}.

\bibitem{chen-volkmer}
Xianghong Chen and Hans Volkmer,
\emph{On transfer operators on the circle with trigonometric weights},
2016. \href{https://arxiv.org/abs/1610.07658}{arXiv:1610.07658}.

\bibitem{bardadyn-kwasniewski}
K.~Bardadyn and B.~Kwa\'sniewski,
\emph{Spectrum of weighted isometries: $C^*$-algebras, transfer operators
and topological pressure}, 2019; revised 2020.
\href{https://arxiv.org/abs/1911.04811}{arXiv:1911.04811}.

\bibitem{baake-coons}
Michael Baake and Michael Coons,
\emph{Correlations of the Thue--Morse sequence}, 2022.
\href{https://arxiv.org/abs/2209.07102}{arXiv:2209.07102}.

\bibitem{ulas-ulas}
Maciej Ulas and Oliwia Ulas,
\emph{On certain arithmetic properties of Stern polynomials}, 2011.
\href{https://arxiv.org/abs/1102.5109}{arXiv:1102.5109}.

\bibitem{arias}
Juan Arias de Reyna,
\emph{An infinitely differentiable function with compact support:
Definition and Properties}, 2017 English translation of the 1982 paper
in \emph{Rev. Real Acad. Ciencias Madrid} \textbf{76}, 21--38.
\href{https://arxiv.org/abs/1702.05442}{arXiv:1702.05442}.

\end{thebibliography}
\end{document}
